\subsection{sel4}
\begin{frame}{Time Protection: Concept}
  \framesubtitle{Time protection: the missing OS abstraction}
	\begin{block}{Definition}
    OS mechanisms preventing interference between security domains that make
    execution speed dependent on other domains' activities.
	\end{block}

	\vspace{0.5cm}

  \pause
	\begin{block}{Two Forms of Partitioning}
		\begin{itemize}
			\item \textbf{Spatial Partitioning}: Separate memory and resources (no sharing)
			\item \textbf{Temporal Partitioning}: Reset state to known defaults (flushing)
		\end{itemize}
	\end{block}
\end{frame}

\begin{frame}{Time Protection: Requirements}
  \framesubtitle{Time protection: the missing OS abstraction}
  \begin{itemize}[<+->]
		\item Flush on-core state on domain switch deterministically
		\item Partition the OS:
		      \begin{itemize}
			      \item Spatially: clone kernel into each domain
			      \item Temporally: reset shared state to known default
			      \item Interrupts: handled deterministically in regards to domains
		      \end{itemize}
		\item Deterministic data sharing for intended channels
	\end{itemize}
\end{frame}

\begin{frame}{Time Protection: Implementation}
  \framesubtitle{Time protection: the missing OS abstraction}
	\begin{block}{Goal}
		Close all possible covert channels

		Argument: closing covert channels closes side channels
	\end{block}
	\begin{itemize}
		\item Touch all l1 cache lines on domain switch
    \item Cache coloring of the l2 cache
    \item Kernel clone implementation, relying on sel4 separation capabilities
    \item Minimize shared kernel state, and guarantee determinism here
	\end{itemize}
\end{frame}
